\section{Experimental Results}
\label{sec:results}

We evaluated the vulnerability of the target architectures under both Black-Box (Genetic Algorithm) and White-Box (Saliency-Guided PGD) threat models. The results quantify the "Sponge Effect" in terms of computational depth, energy consumption, and transferability.

\subsection{Attack Efficacy: Black-Box vs. White-Box}
A key contribution of this work is the comparative analysis between derivative-free optimization and gradient-based attacks. Figure \ref{fig:ponder_vis} (Top Right) illustrates the relative efficacy of each approach.

\subsubsection{Computational Depth}
The ACT-LSTM baseline processes benign inputs with an average computational depth of $\mu=5.48$ ponder steps per time-step.
\begin{itemize}
    \item \textbf{Genetic Algorithm (GA):} The black-box GA latency attack increased this average to $9.30$ steps ($+70\%$), directly optimizing wall-clock latency without gradient access.
    \item \textbf{PGD:} The gradient-based PGD attack achieved an average of $6.38$ steps ($+16\%$), optimizing the differentiable ponder loss proxy.
\end{itemize}
Contrary to the common assumption that white-box access yields strictly superior attacks, the GA outperforms PGD on the ACT-LSTM. As shown in Figure \ref{fig:ponder_vis} (Top Left), the GA latency attack (purple) forces the model to hit the hard cap of 20 ponder steps at multiple intervals, a behavior less frequent under PGD. The histogram (Bottom Right) confirms this, showing a heavier tail for GA attacks. This reversal occurs because the GA directly maximizes measured latency, whereas PGD optimizes a differentiable proxy ($\mathcal{L}_{ponder}$) that does not perfectly correlate with wall-clock time.

% --- TABLE II: COMPARISON ---
\begin{table}[t]
    \centering
    \caption{Comparison of Attack Efficacy on ACT-LSTM. The Black-Box GA achieves greater ponder depth and latency than White-Box PGD by directly optimizing wall-clock time rather than a differentiable proxy.}
    \label{tab:attack_comparison}
    \renewcommand{\arraystretch}{1.2}
    \resizebox{\columnwidth}{!}{%
    \begin{tabular}{l c c c c}
        \toprule
        \textbf{Method} & \textbf{Ponder Depth} & \textbf{Latency (ms)} & \textbf{Energy (mJ)}\\
        \midrule
        \textbf{Baseline (Benign)} & 5.48 & 47.99 & 818.44\\
        \midrule
        \textbf{Black-Box (GA)} & \textbf{9.30} (\textbf{+70\%}) & $\approx$\textbf{73.5} (\textbf{+53\%})  & $\approx 1785$ (\textbf{+78\%}) \\
        \textbf{White-Box (PGD)} & 6.38 (\textbf{+16\%}) & 55.74 (\textbf{+16\%}) & 1023.74 (\textbf{+25\%}) \\
        \bottomrule
    \end{tabular}%
    }
\end{table}

\begin{figure*}[t]
    \centering
    \includegraphics[width=\textwidth]{act_ponder_visualization.png}
    \caption{\textbf{Dynamics of the Sponge Attack.} (Top Left) Ponder steps per time-increment under each attack variant. (Top Right) Average ponder steps: the GA latency attack achieves the highest depth (9.30), followed by GA energy (8.97), PGD energy (6.81), and PGD latency (6.38), all above the 5.48 baseline. (Bottom Right) Histogram showing that GA attacks push more time-steps into the high-ponder regime.}
    \label{fig:ponder_vis}
\end{figure*}

\subsection{Energy and Latency Impact}
The increase in computational depth translates directly to physical resource exhaustion.
\begin{itemize}
    \item \textbf{Latency:} Under the PGD latency attack, the per-sample latency increased from $47.99$ ms to $55.74$ ms ($+16\%$). The GA latency attack achieved a larger increase to $\approx 73.5$ ms ($+53\%$) (see Fig. \ref{fig:latency_res} in the Appendix).
    \item \textbf{Energy:} The total energy consumption per inference increased from $818.44$ mJ to $1023.74$ mJ ($+25\%$) under PGD. The GA energy attack achieved $\approx 1785$ mJ ($+78\%$) (see Fig. \ref{fig:energy_res} in the Appendix).
\end{itemize}
Crucially, the Power draw remained constant at $\approx 17$ W under PGD. This confirms that the "Sponge" effect is purely temporal; the attack does not increase the instantaneous intensity of the computation but rather prolongs its duration, draining battery-constrained devices over time.

\subsection{XAI Analysis: Availability vs. Accuracy}
To understand the mechanism of the attack, we analyzed the feature attributions using Integrated Gradients (IG). Figure \ref{fig:xai_analysis} reveals a critical distinction between features important for \textit{accuracy} and those important for \textit{availability}.

The "Feature Importance" chart (Bottom Right) shows that while the \textbf{Power} feature has the highest IG Attribution (blue bar)—meaning it is critical for the prediction—the attack focuses its perturbation energy (orange bars) on environmental features like \textbf{Temperature}, \textbf{Wind Direction}, and \textbf{Wind Gusts}.
This suggests that while the model looks at previous power output to predict the next value, it relies on complex environmental variables to determine "how hard to think." The attacker exploits this by creating "chaotic weather" scenarios (high gusts/variable temperature) that maximize internal uncertainty.

\begin{figure}[h]
    \centering
    \includegraphics[width=\columnwidth]{act_pgd_latency_xai_analysis.png}
    \caption{\textbf{XAI-Guided Attack Analysis.} (Bottom Right) A key insight: While 'Power' (blue bar) is the most important feature for prediction, the attacker targets 'Wind Gusts' and 'Temperature' (orange bars) to cause latency. The model 'thinks' hardest during chaotic weather.}
    \label{fig:xai_analysis}
\end{figure}

\subsection{Attack Transferability and Defense}
Finally, we evaluated whether sponge examples generated for the ACT-LSTM affect static architectures. Figure \ref{fig:transferability} shows the cross-model performance.
\begin{itemize}
    \item \textbf{ACT-LSTM:} Sufferers massive spikes (Blue bars) under its own adversarial examples.
    \item \textbf{DeepAR (Static):} Remains completely unaffected (Orange bars), with latency staying flat even when fed the ACT-optimized adversarial inputs.
\end{itemize}
This confirms that \textbf{latency-based sponge attacks are architecture-specific}. The halting-unit loops exploited in the ACT model do not exist in the static DeepAR graph. Architectural distillation (replacing dynamic models with static ones) is therefore an effective defense against \textit{latency} attacks—though not against the bit-flip energy attack, which succeeds on all architectures (Section~\ref{sec:bitflip_results}).

\begin{figure}[h]
    \centering
    \includegraphics[width=\columnwidth]{attack_transferability_analysis.png}
    \caption{\textbf{Latency Transferability Analysis.} Latency attacks optimized for the dynamic ACT-LSTM (Blue) cause significant spikes on the target but have zero effect on the static DeepAR model (Orange), confirming that static architectures are robust to latency-based sponge attacks.}
    \label{fig:transferability}
\end{figure}

\subsection{Energy via Switching Activity: Bit-Flip Analysis}
\label{sec:bitflip_results}
While the preceding attacks increase energy by prolonging computation, the Bit-Flip Oracle targets a complementary mechanism: maximizing dynamic power dissipation by increasing transistor switching activity. We evaluated this attack across all three architectures. Table~\ref{tab:bitflip_comparison} summarizes the results.

% --- TABLE III: BIT-FLIP COMPARISON ---
\begin{table}[t]
    \centering
    \caption{Bit-Flip Oracle Results Across Architectures. Flip counts scale with first-layer weight count; the GA achieves +25--36\% increase across all models.}
    \label{tab:bitflip_comparison}
    \renewcommand{\arraystretch}{1.2}
    \resizebox{\columnwidth}{!}{%
    \begin{tabular}{l r r r r}
        \toprule
        \textbf{Model} & \textbf{First-Layer Weights} & \textbf{Baseline Flips} & \textbf{Attack Flips} & \textbf{Increase} \\
        \midrule
        ACT-LSTM & 4,608 & 64,859 & 82,237 & \textbf{+26.8\%} \\
        DeepAR & 18,432 & 261,690 & 327,730 & \textbf{+25.2\%} \\
        Chronos-v2 & 2,097,152 & 28,387,055 & 38,749,216 & \textbf{+36.5\%} \\
        \bottomrule
    \end{tabular}%
    }
\end{table}

The key finding is the strong proportionality between the number of first-layer weights and the absolute flip count. Chronos-v2, with its 2M-parameter embedding layer (\texttt{shared.weight}, $4096 \times 512$), produces approximately $470\times$ more bit-flip transitions than ACT-LSTM (4,608 weights) and $118\times$ more than DeepAR (18,432 weights) under adversarial input.

Figure~\ref{fig:bitflip_comparison} summarizes the results across all three models. The left panel (log scale) shows the absolute flip counts, illustrating how Chronos-v2's large embedding layer produces orders of magnitude more switching transitions than the smaller models. The right panel shows the relative attack efficacy: the GA achieves a consistent +25--36\% increase in switching activity across all architectures. Per-model GA optimization trajectories are provided in the Appendix (Figs.~\ref{fig:act_bitflip}--\ref{fig:chronos_bitflip}).

\begin{figure}[h]
    \centering
    \includegraphics[width=\columnwidth]{bitflip_comparison.png}
    \caption{\textbf{Bit-Flip Oracle: Cross-Model Comparison.} (Left) Absolute switching transitions on log scale—Chronos-v2 produces $\sim$470$\times$ more flips than ACT-LSTM due to its larger first layer. (Right) Relative increase: the GA achieves +25--36\% across all architectures, confirming the attack is architecture-agnostic.}
    \label{fig:bitflip_comparison}
\end{figure}

Notably, unlike the latency-based sponge attacks which are effective primarily against dynamic architectures (Section~\ref{sec:discussion}), the bit-flip attack succeeds across \textit{all} architectures—including the static DeepAR and the quantized Chronos-v2. This is because switching activity depends on the data representation flowing through the first layer, not on the computational graph's dynamism. The bit-flip vector thus provides an architecture-agnostic path to increasing energy consumption.
